\uuid{gpn0}
\titre{DL usuels/Calcul de limites}
\auteur{hanani ; mzari}
\organisation{exo7}
\licence{CC BY-NC-SA 4.0}
\datecreate{2020-01-20}
\niveau{L1}
\module{Analyse}
\chapitre{Développement limité}
\sousChapitre{Applications}
\difficulte{2}
\theme{Applications des DL, DL usuels, Calcul de limites}
\exo7qcmid{lille-2.developpements-limites.032}

\contenu{
\texte{
Soit $a\in \Rr$ et $f$ une fonction. On dit que $f(x)$ est un $o\left((x-a)^n\right)$ au voisinage de $a$ et on écrit $f(x)=o((x-a)^n)$ si $\displaystyle \lim _{x\to a}\frac{f(x)}{(x-a)^n}=0$. On note par $DL_n(a)f(x)$ le développement limité de $f$ à l'ordre $n$ en $a$. On note aussi $DL_n(a^+)f(x)$ (resp. $DL_n(a^-)f(x)$) le développement limité de $f$ à l'ordre $n$ à droite en $a$ (resp. à gauche en $a$).
}

\begin{question}
\qclassification{125.03}
\qtags{motcle=DL usuels/Calcul de limites}
Parmi les égalités suivantes, cocher celles qui sont vraies :
\begin{answers}
\bad{$\displaystyle \lim _{x\to 0}\left(1+x\right)^{1/x}=1$.}
\good{$\displaystyle\lim _{x\to 0}\left(\cos x\right)^{1/\sin ^2x}=\frac{1}{\sqrt{\mathrm{e}}}$.}
\good{$\displaystyle \lim _{x\to 0}\left(\frac{1}{\sin ^2x}-\frac{1}{x^2}\right)=\frac{1}{3}$.}
\bad{$\displaystyle \lim _{x\to 0}\frac{\cos x-\sqrt{1-x^2}}{x^2\sin ^2x}$ n'existe pas.}
\end{answers}
\vskip2mm
\begin{explanations}
On a : $\displaystyle \left(1+x\right)^{1/x}=\mathrm{e}^{\frac{\ln (1+x)}{x}}$ et $\displaystyle \frac{\ln (1+x)}{x}=1+o(1)$. Donc
$$\mathrm{e}^{\frac{\ln (1+x)}{x}}=\mathrm{e}+o(1) \Rightarrow \lim _{x\to 0}\left(1+x\right)^{1/x}=\mathrm{e}.$$
De même, $\displaystyle \left(\cos x\right)^{1/\sin ^2x}=\mathrm{e}^{\ln (\cos x)/\sin ^2x}$. Or $\sin ^2x=x^2+o(x^2)$ et
$$\displaystyle \ln (\cos x)=\ln \left(1-\frac{x^2}{2}+o(x^2)\right)=-\frac{x^2}{2}+o(x^2).$$
Donc
$$\frac{\ln (\cos x)}{\sin ^2x}=-\frac{1}{2}+o(1)\Rightarrow \lim _{x\to 0}\left(\cos x\right)^{1/\sin ^2x}=\mathrm{e}^{-1/2}.$$
Ensuite, $\displaystyle \frac{1}{\sin ^2x}-\frac{1}{x^2}=\frac{x^2-\sin ^2x}{x^2\sin ^2x}$, et le $\displaystyle DL_4(0)\sin ^2x=x^2-\frac{x^4}{3}+o(x^4)$ donne
$$DL_0(0)\frac{x^2-\sin ^2x}{x^2\sin ^2x}=\frac{1}{3}+o(1)\Rightarrow \lim _{x\to 0}\left(\frac{1}{\sin ^2x}-\frac{1}{x^2}\right)=\frac{1}{3}.$$
Enfin, $\displaystyle \cos x=1-\frac{x^2}{2}+\frac{x^4}{24}+o(x^4)$, $\displaystyle \sqrt{1-x^2}=1-\frac{x^2}{2}-\frac{x^4}{8}+o(x^4)$ et $x^2\sin ^2x=x^4+o(x^4)$. Donc
$$\frac{\cos x-\sqrt{1-x^2}}{x^2\sin ^2x}=\frac{x^4/6+o(x^4)}{x^4+o(x^4)}=\frac{1}{6}+o(1)\Rightarrow \lim _{x\to 0}\frac{\cos x-\sqrt{1-x^2}}{x^2\sin ^2x}=\frac{1}{6}.$$
\end{explanations}
\end{question}
}
